% !TeX root = ../抽象代数笔记.tex
% 由主文档通过 \input 载入；请编译主文档。
\chapter{绪论与全书大纲}

\section{研究对象与课程定位}

整数可以相加、相乘，几何变换可以复合，矩阵可以相乘，多项式可以作带余除法。这些运算的具体对象不同，却反复呈现结合律、单位元、逆元和分配律等共同规律。\emph{抽象代数（abstract algebra）}从这些规律出发，研究带有运算的集合，以及保持运算的映射。

本书围绕三类问题组织内容：一个代数结构由哪些元素生成；哪些映射能够保持它的运算；怎样借助子结构、商结构和分解定理理解整体。\emph{群（group）}用于研究单一可逆运算与对称性，\emph{环（ring）}和\emph{域（field）}用于研究加法、乘法与方程，\emph{模（module）}则把线性空间中的标量推广到环，连接群论、环论和线性代数。

课程按本科两学期的完整学习范围设计，预备知识为集合与映射、整数整除和基础线性代数。必要的证明方法与整数知识在正文中补足。群、环、域构成核心主线；模论、标准形及标为选学的专题用于深化。微积分不是主要先修课程，研究生层次的交换代数、表示论和同调代数不在本书系统展开的范围内。

本书按上述路线逐章展开。第\ref{ch:preliminaries}至\ref{ch:cyclic-permutation}章已写成正文，给出定义、证明与例题；其余章节暂保留写作大纲，以说明后续内容及依赖关系。大纲中列出的证明任务不视为已经完成的论证。

\section{全书课程大纲}

\begin{description}
  \item[第一篇：预备知识与群论] 从数学语言和对称性出发，建立群的基本工具，再研究有限群的结构。
\end{description}
\begin{enumerate}
  \item 第\ref{ch:preliminaries}章：数学语言与预备知识。
  \item 第\ref{ch:groups}章：群、子群与生成。
  \item 第\ref{ch:cyclic-permutation}章：循环群与置换群。
  \item 第\ref{ch:cosets}章：陪集、指数与拉格朗日定理。
  \item 第\ref{ch:quotient-groups}章：正规子群、商群与同态定理。
  \item 第\ref{ch:actions}章：群作用与计数。
  \item 第\ref{ch:sylow}章：西罗定理与有限群。
  \item 第\ref{ch:abelian}章：直积与有限阿贝尔群。
  \item 第\ref{ch:solvable-groups}章：群列、单群与可解群。
\end{enumerate}

\begin{description}
  \item[第二篇：环与多项式] 把整数与多项式中的整除规律放入统一框架，以理想构造商环，并为扩域作准备。
\end{description}
\begin{enumerate}[resume]
  \item 第\ref{ch:rings}章：环、整环、域与环同态。
  \item 第\ref{ch:ideals}章：理想、商环与中国剩余定理。
  \item 第\ref{ch:factorization}章：整除性与唯一分解。
  \item 第\ref{ch:polynomials}章：多项式环与不可约性。
\end{enumerate}

\begin{description}
  \item[第三篇：域与伽罗瓦理论] 从添加方程的根出发，用域扩张描述代数数，再用根的对称性研究方程能否用根式求解。
\end{description}
\begin{enumerate}[resume]
  \item 第\ref{ch:extensions}章：域扩张与代数元素。
  \item 第\ref{ch:constructions}章：尺规作图与二次扩张。
  \item 第\ref{ch:splitting}章：分裂域、可分扩张与正规扩张。
  \item 第\ref{ch:finite-fields}章：有限域。
  \item 第\ref{ch:galois}章：伽罗瓦对应与伽罗瓦群计算。
  \item 第\ref{ch:radicals}章：根式扩张与方程的可解性。
\end{enumerate}

\begin{description}
  \item[第四篇：模与线性代数的统一] 把阿贝尔群和线性变换放入模的语言，利用结构定理解释标准形。
\end{description}
\begin{enumerate}[resume]
  \item 第\ref{ch:modules}章：模、子模与模同态。
  \item 第\ref{ch:pid-modules}章：主理想整环上的有限生成模与标准形。
\end{enumerate}

\section{学习路线与知识依赖}

初学时先完成第\ref{ch:preliminaries}至\ref{ch:actions}章，掌握定义、例子、同态和商结构的基本方法，再学习西罗定理与有限群分解。环论从第\ref{ch:rings}章开始，主要依赖前面的群论基础；第\ref{ch:sylow}至\ref{ch:solvable-groups}章的较深内容可以与环论交替学习。

域论依赖第\ref{ch:polynomials}章的不可约多项式与商环构造，以及线性代数中的基、维数和线性无关。第\ref{ch:constructions}章是域扩张的独立应用，略过它不影响后续理论。第\ref{ch:radicals}章同时依赖可解群与伽罗瓦对应。模论在完成环论后即可开始，不以伽罗瓦理论为先修。

一学期入门路线可以选取第\ref{ch:preliminaries}至\ref{ch:quotient-groups}章、群作用的基本部分、第\ref{ch:rings}至\ref{ch:extensions}章，并以有限域或尺规作图作结。两学期路线覆盖完整核心主线，再按学习进度加入模论与选学内容。具体进度应以能独立完成定义验证、计算与基础证明为准。

\section{统一记号与学习要求}

整数、有理数、实数、复数分别记为 $\mathbb Z,\mathbb Q,\mathbb R,\mathbb C$；自然数约定为 $\mathbb N=\{0,1,2,\ldots\}$。群通常记为 $G$，单位元记为 $e$，子群关系记为 $H\leq G$；加法群的单位元记为 $0$。整数 $n\geq 1$ 的剩余类集合记为 $\mathbb Z/n\mathbb Z$，整数 $a$ 的剩余类记为 $\overline a$。映射复合 $f\circ g$ 表示先作用 $g$，再作用 $f$。

环均指有乘法单位元的结合环，允许零环；环同态与子环均要求保持单位元。整环和域额外要求 $1\ne0$。环不默认交换，进入交换环的结论时会明确说明。模默认指含幺左模。\emph{同态（homomorphism）}保持相应运算，\emph{同构（isomorphism）}则是可逆的结构保持映射；同构记为 $\cong$，不与集合相等混用。

每章展开时采用“问题与例子、定义、基本性质、关键定理及证明、应用、习题”的顺序。习题包含定义辨析、计算、证明和综合应用；反例用于说明假设不可随意删除。正文中正式引入符号与术语，各章大纲中的“练习与自检”说明后续题组应检验的能力。

